% ---------------------------------------------------------------- 摘要（第二部分分册）
% 与 tex/abstract.tex 同源，仅供“第二部分”分册使用：
%   * 保留全报告定位一句与“第二部分的主要发现”六条；
%   * 删去“第一部分的主要发现”四条（其口径、样本与结论属于第一部分，见另册）；
%   * 本册保留的 \ref 全部指向第二部分自身的章节（第 7 章、第 8--13 章），
%     分册中这些标签均已定义，不会印成 ??。
\section*{摘要}
\addcontentsline{toc}{section}{摘要}

\begin{keybox}[报告要点]
\small
本报告以\textbf{申万行业分类（2021 版）}为统一框架，分为两大部分：
第一部分从供需关系出发分析农业各子行业的周期性，以\hl{猪周期为基准}
测算跨品种、跨行业的相关性与传导时滞；第二部分系统梳理
申万农林牧渔行业\textbf{全部 \num{104} 家 A 股上市公司}近三年的经营与融资画像。

\textbf{本册为第二部分}，对应第七章至第十五章（逐公司画像、融资历史与潜在融资需求）
以及附录 A、附录 B；第一部分（子行业周期性与跨品种相关性）为另一册。

\textbf{第二部分的主要发现}：
（1）\hl{第二部分统一采用 2026 年半年报口径（合并报表、未年化），
市值一律采用总市值口径}（2026 年 10 月 12 日收盘价 $\times$ 总股本，
不使用流通市值）：100 余家公司的总市值合计 \num{12237} 亿元，
前三家占 \num{38.2}\%；
\num{104} 家公司中 \num{46} 家亏损（占 \num{44.2}\%）、\num{44} 家半年 ROE 为负，
ROE 中位数仅 \num{0.75}\%，半年 ROE 超过 \num{5}\%（约合年化 \num{10}\%）的只有 \num{16} 家；
其中生猪养殖 \num{11} 家在 2026 年上半年\hl{全部亏损}、合计 \num{-163.8} 亿元
（上年同期为 \num{+154.1} 亿元）；
（2）子行业分化鲜明：动保（半年 ROE 中位数 \num{2.95}\%、亏损面 \num{11.1}\%、
资产负债率中位数 \num{30.2}\%）与海洋捕捞（\num{9.79}\%）为抗周期方向，
饲料（半年 ROE 中位数 \num{-6.17}\%、亏损面 \num{66.7}\%）为最受损环节；
（3）\num{14} 家公司同时处于 2026 年上半年亏损与 6 月末资产负债率 $>$ \num{65}\% 状态
（其中 \num{6} 家 $>$ \num{80}\%），\num{31} 家连续三年未分红；
据财务约束与公告行为推断，第一档“补血型”融资需求按最低口径为 \num{150}---\num{250} 亿元、
按“负债率压回 \num{65}\%”的重组口径约 \num{305} 亿元
（若改以偿还债务达标则需减少负债约 \num{567} 亿元），
\hl{当前农业上市公司的资本运作主题是修复资产负债表，而非扩大产能}；
（4）与 2025 年年报（\num{34} 家亏损、ROE 中位数 \num{2.64}\%）相比，
\hl{2026 年半年报的亏损面进一步扩大}：归母净利润合计由上年同期的 \num{+270.9} 亿元
转为 \num{-126.1} 亿元，而营业收入合计 \num{6085} 亿元并未同幅萎缩，
\hl{报表层面的底部晚于价格层面的底部}；
（5）逐公司画像给出三处关键事实：\hl{北大荒} 2026 年上半年由盈转亏，
主因是补缴 2021---2025 年企业所得税与滞纳金 \num{14.10} 亿元（一次性计入当期损益、
不追溯调整），剔除后扣非亏损仅 \num{2.09} 亿元，土地发包主业并未恶化；
\hl{种业去产能已经开始}，全国玉米制种面积自 2023 年的 \num{462} 万亩连续三年降至
2026 年的 \num{358.5} 万亩（\num{-24}\%），但 \num{175}\% 的供需比意味着库存出清尚需时间；
\hl{诺普信}的利润已主要来自基质栽培、错季上市的蓝莓（子公司爱莓庄上半年净利润 \num{6.07} 亿元），
农药制剂业务同期仍稳定增长 \num{2.4}\%；
（6）在周期底部真正区分公司的是\hl{\emph{在低点而非高点加杠杆扩张的能力}}——
在景气高点借来的产能会把成本锁在最高处，只有在行业普遍亏损时仍能拿到低成本资金的公司
才有条件承接别人的产能。
第二部分的各行业章（第 \ref{sec:planting} 章至第 \ref{sec:fishery} 章）末尾设有“上市公司画像”小节，
给出全部 \num{104} 家公司的单独画像。画像\hl{逐家检索公开资料}写成，
不使用统一模板：每家公司的沿革与股权变动、融资历史、风险事件与战略动作各不相同，
正文按“基本情况、股权结构与控制权、经营要点、所处行业的基本面（供需）、
与同类公司比较的竞争优势、股价与周期、盈利与财务、战略转型与融资”的顺序展开，
每条事实均可追溯至公司公告、定期报告或公开报道（数据口径见第 \ref{sec:comp-overview} 章）；
每家配股价与行业指数对比图、近年主要财务指标图
（营业收入与归母净利润、ROE 与资产负债率、货币资金与有息负债）、
关键财务指标表与历史融资/资本运作表。
对经历过重整或债务重组的公司（如天邦食品、正邦科技、傲农生物、华英农业等）
单独梳理其融资与重组过程，并按最新财务指标判断未来的潜在融资需求
（以短期债务缺口与年化经营缺口两个口径度量）。
\end{keybox}

\vspace{6pt}
\noindent\textbf{关键词}\quad 农业周期；猪周期；跨品种相关性；申万行业分类；
上市公司画像；融资需求
\clearpage
